\begin{abstract}
Blink detection from electroencephalography (EEG) serves two complementary purposes:
blinks are a major physiological contaminant that must be identified for artifact
handling, and they are also informative behavioral markers linked to arousal, fatigue,
cognitive load, and operational performance.
Despite substantial progress in blink detection, a practically important setting
remains unaddressed: epoch-structured recordings in which a long continuous signal is
divided into non-overlapping analysis windows and some windows are pre-rejected by the
analyst before automated processing begins.
This setting is the norm in driving neuroscience and other neuroergonomic paradigms,
yet existing detectors---including BLINKER and MNE-based amplitude routines---operate
on concatenated signals and implicitly pool blink-containing and blink-free epochs when
estimating thresholds.
We show that this pooling systematically biases the threshold upward and reduces recall.

We propose a three-stage epoch-aware pipeline.
Stage~A screens epochs by their peak-to-peak amplitude via an autoreject-inspired
cross-validation criterion, yielding a set of suspicious (likely blink-containing)
epochs.
Stage~B estimates a robust sample-level threshold from the suspicious epoch pool using
the median and median absolute deviation (MAD), concentrating estimation on
signal segments enriched for blink-like excursions.
Stage~C applies the threshold to extract blink regions from the suspicious epochs and
maps detections back to epoch-local coordinates.
A fallback mechanism ensures stability when too few suspicious epochs are available.

We evaluate five conditions---BLINKER with naive concatenation, direct Bayesian
threshold optimisation, MNE amplitude annotation, and the proposed pipeline with
mean-based and median-based estimators---on a large naturalistic driving corpus
under epoch durations of 20, 30, 60, and 120 seconds.
Statistical comparisons use macro-averaged event-level F1 with Wilcoxon signed-rank
tests.
Additional analyses examine sensitivity to the minimum suspicious-epoch count
$n_{\min}$, stability across boundary tolerances, robustness under user-initiated
epoch pre-rejection, and morphological characteristics of missed and spurious blinks.
Results demonstrate that epoch-aware threshold estimation consistently outperforms
naive concatenation, that the robust MAD-based estimator is preferable to a
mean-based alternative, and that the pipeline degrades gracefully as the available
epoch pool shrinks due to pre-rejection.
\end{abstract}
